\documentclass[11pt]{article}
\usepackage[T1]{fontenc}
\usepackage[utf8]{inputenc}
\usepackage{lmodern,amsmath,amssymb,amsthm,booktabs,array}
\usepackage[margin=1in]{geometry}
\usepackage[hidelinks]{hyperref}
\newtheorem{theorem}{Theorem}[section]
\newtheorem{lemma}[theorem]{Lemma}
\newtheorem{proposition}[theorem]{Proposition}
\newtheorem{corollary}[theorem]{Corollary}
\newtheorem{example}[theorem]{Example}
\theoremstyle{remark}\newtheorem{remark}[theorem]{Remark}
\newcommand{\Fp}{\mathbf F_p}
\newcommand{\Fthree}{\mathbf F_3}
\newcommand{\Cay}{\operatorname{Cay}}
\newcommand{\Hom}{\operatorname{Hom}}
\newcommand{\Spec}{\operatorname{Spec}}
\newcommand{\Exp}{\operatorname{exp}}
\newcommand{\Z}{\mathbf Z}
\newcommand{\Q}{\mathbf Q}
\newcommand{\T}{\mathcal T}
\title{Cospectrality of a structured three-element support\newline under prime cyclic translation}
\author{Alexandre Couret\\Independent researcher, France}
\date{4 October 2026\\\small Version 0.1 for independent mathematical review}
\begin{document}
\maketitle
\begin{abstract}
Let $H$ be a finite abelian group generated by distinct elements $\kappa,\kappa'$, and let $p$ be an odd prime. In $H\times\mathbf F_p$, consider the support $T=\{(\kappa,0),(\kappa,1),(\kappa',0)\}$ and its translate $T^+=T+(0,1)$. For $p\geq5$, the associated Cayley digraphs are cospectral exactly when there is a homomorphism $h:H\to\mathbf F_p$ with $h(\kappa')=1$ and $h(\kappa)\in\{1,2\}$. This condition yields an explicit group automorphism between the two supports. For $p=3$, the same criterion applies when $3$ divides $|H|$. When $3$ does not divide $|H|$, cospectrality is equivalent to invariance of the character-value subgroup $\Gamma=\{(\psi(\kappa),\psi(\kappa')):\psi\in\widehat H\}$ under $(u,v)\mapsto(-v,-u)$. The arguments use short vanishing sums of roots of unity for $p\geq5$ and a unit-circle chord argument for $p=3$. We also give formulations in terms of the relation lattice and identify a previously studied circulant example within this family.
\end{abstract}
\noindent\textbf{Keywords.} Cayley digraph; cospectrality; roots of unity; finite abelian group; circulant.\\
\textbf{MSC 2020.} 05C50; 05C25; 05C60.

\section{Introduction and scope}
Cospectrality does not in general determine a graph up to isomorphism. For circulants, this question has been studied through arithmetic conditions on the support, the spectral \`Ad\'am property, and constructions of cospectral nonisomorphic pairs; see Brown~\cite{Brown}, Mans, Pappalardi and Shparlinski~\cite{MPS}, and M\"onius~\cite{Moenius2020,Moenius}. The survey of Liu and Zhou~\cite{LZ} provides further references. Here we consider a prescribed pair of three-element supports rather than arbitrary supports of size three.

All groups are written additively. Let $H=\langle\kappa,\kappa'\rangle$ be finite abelian, with $\kappa\ne\kappa'$. We identify the chosen generator of $C_p$ with $1\in\Fp$ and put
\begin{equation}\label{supports}
G=H\times\Fp,\qquad
T=\{(\kappa,0),(\kappa,1),(\kappa',0)\},\qquad T^+=T+(0,1).
\end{equation}
Our convention is that $\Cay(G,S)$ has an arc $g\to g+s$ for each $s\in S$. A loop is retained if $0\in S$. The adjacency operator acts by $(Af)(g)=\sum_{s\in S}f(g+s)$. Cospectrality means equality of the multisets of complex adjacency eigenvalues, including algebraic multiplicities. No symmetrization of the support is made.

The classification below is limited to the pair~\eqref{supports}. In particular, it asserts neither a general CI property nor determination by spectrum for all abelian Cayley digraphs of out-degree three. Previous cospectral examples are included rather than claimed anew. The present version consolidates the arguments of the dated source notes 10A6, 10C1, 10C2 and 10C3. Independent checking of these arguments and a complete assessment of priority remain pending; no claim of priority is made in this version.

\section{Classification}
Let $\widehat H$ be the complex character group of $H$, and define
\begin{equation}\label{Gamma}
\Gamma=\{(\psi(\kappa),\psi(\kappa')):\psi\in\widehat H\}.
\end{equation}
It is a finite subgroup of the product of two unit circles. Since $\kappa,\kappa'$ generate $H$, evaluation is injective and $|\Gamma|=|H|$.

\begin{theorem}\label{main5}
For $p\geq5$, the following are equivalent:
\begin{enumerate}
\item $\Cay(G,T)$ and $\Cay(G,T^+)$ are cospectral;
\item there exists $h\in\Hom(H,\Fp)$ such that $h(\kappa')=1$ and $h(\kappa)\in\{1,2\}$;
\item a group automorphism of $G$ sends $T$ onto $T^+$;
\item the two Cayley digraphs are isomorphic.
\end{enumerate}
\end{theorem}

\begin{theorem}\label{main3}
Let $p=3$ and $\beta(u,v)=(-v,-u)$.
\begin{enumerate}
\item If $3\nmid|H|$, the two digraphs are cospectral if and only if $\beta(\Gamma)=\Gamma$.
\item If $3\mid|H|$, the four conditions in Theorem~\ref{main5}, with $\Fp=\Fthree$, are equivalent.
\end{enumerate}
For all $H$, the condition $\beta(\Gamma)=\Gamma$ is sufficient for cospectrality; when $3\mid|H|$, it implies condition~2 above.
\end{theorem}

\begin{corollary}\label{combined}
For $p=3$, cospectrality is equivalent to the disjunction of the homomorphism condition and $\beta(\Gamma)=\Gamma$.
\end{corollary}
Indeed, the homomorphism condition is impossible if $3\nmid|H|$, and $\beta(\Gamma)=\Gamma$ implies it in the complementary regime.

\begin{lemma}\label{spectral}
For $(u,v)\in\Gamma$ and $z\in\mu_p$, the eigenvalues, counted once for each such pair and phase, are
\[
\lambda_T(u,v,z)=u(1+z)+v,\qquad
\lambda_{T^+}(u,v,z)=z u+z^2u+zv=z\lambda_T(u,v,z).
\]
\end{lemma}
\begin{proof}
The characters $f(x,y)=\psi(x)z^y$ form an eigenbasis of the adjacency operator. Substituting the three elements of each support proves the displayed formulas, with the required multiplicities.
\end{proof}

\begin{lemma}\label{suffh}
For every odd prime $p$, condition~2 of Theorem~\ref{main5} gives a group automorphism carrying $T$ onto $T^+$.
\end{lemma}
\begin{proof}
If $h(\kappa)=1$, set $\phi_+(x,y)=(x,y+h(x))$. If $h(\kappa)=2$, set $\phi_-(x,y)=(x,h(x)-y)$. The first map has inverse $(x,y)\mapsto(x,y-h(x))$ and the second is an involution. The images of the three support elements are respectively
\[
\begin{array}{c|ccc}
 & (\kappa,0)&(\kappa,1)&(\kappa',0)\\\hline
\phi_+ &(\kappa,1)&(\kappa,2)&(\kappa',1)\\
\phi_- &(\kappa,2)&(\kappa,1)&(\kappa',1).
\end{array}
\]
Thus $\phi_\pm(T)=T^+$. Group automorphisms give digraph isomorphisms, which imply cospectrality.
\end{proof}

\begin{lemma}\label{dualh}
Fix a primitive $p$th root $z$. The homomorphism condition is equivalent to $\Gamma$ containing at least one of
\[
(z^{-1},z^{-1}),\qquad (z^2,z).
\]
\end{lemma}
\begin{proof}
If $h$ exists, the character $x\mapsto z^{h(x)}$ has pair $(z,z)$ or $(z^2,z)$; inversion gives the first listed pair in the first case. Conversely, a character with either listed pair takes all its values in $\mu_p$, because its values on generators do. Writing it as $z^{h_0(x)}$ produces a homomorphism into $\Fp$. In the first case replace $h_0$ by $-h_0$; the second case has the required values directly.
\end{proof}

\section{Necessity for primes at least five}
We use the following consequence of the classification in Poonen and Rubinstein~\cite[Section~3, Theorem~3 and Table~1]{PR}. A vanishing sum is treated as a multiset of roots, so repetitions count toward its weight.
\begin{lemma}[Short root relations]\label{roots}
Up to multiplying all terms by one root of unity, the minimal vanishing sums of weight at most six have weights two, three, five or six. The first three are regular root polygons $R_2,R_3,R_5$. The weight-six type has six distinct terms and a representative
\[
\zeta_6+\zeta_6^{-1}+\zeta_5+\zeta_5^2+\zeta_5^3+\zeta_5^4=0.
\]
Every vanishing multiset decomposes into minimal vanishing submultisets.
\end{lemma}
The last statement follows by successively removing a minimal vanishing submultiset. We do not claim this root classification as a contribution of the present work. General constraints on possible weights of vanishing sums are studied by Lam and Leung~\cite{LL}; here the explicit weight-six type is needed.

\begin{lemma}\label{forced}
If $p\geq5$, $z$ is primitive of order $p$, and unit roots $A,B,C$ satisfy $A+B+C=2+z$, then $\{A,B,C\}=\{1,1,z\}$ as multisets. If $A+B+C=1+2z$, then $\{A,B,C\}=\{1,z,z\}$.
\end{lemma}
\begin{proof}
For the first assertion, consider $A+B+C-1-1-z=0$. The two occurrences of $-1$ exclude the minimal weight-six type. Lemma~\ref{roots} therefore gives either three $R_2$ blocks or two $R_3$ blocks: a block of weight five would leave one term, and there is no minimal block of weight one or four. In the two-$R_3$ case the two $-1$ occurrences belong to different blocks. The occurrence $-z$ belongs to one of those blocks, so its quotient with $-1$ is a primitive cube root. This contradicts $p\geq5$. In the three-$R_2$ case none of $-1,-1,-z$ can pair with another of these three terms, since none are opposite. Each must pair with its positive opposite. The second assertion is proved in the same way, using the two occurrences of $-z$.
\end{proof}

\begin{proof}[Proof of Theorem~\ref{main5}]
Only necessity remains. Let $z$ be primitive of order $p$. The trivial character of $H$ and phase $z$ give $2+z$ in the spectrum of $T$. Cospectrality and Lemma~\ref{spectral} yield a pair $(u,v)\in\Gamma$ and $z'\in\mu_p$ with
\[
A+B+C=2+z,\qquad A=z'u,\quad B=z'^2u,\quad C=z'v.
\]
By Lemma~\ref{forced}, the following table lists all placements, taking account of $B=z'A$:
\[
\begin{array}{c|c|c|c|c}
A&B&C&z'&(u,v)\\\hline
1&z&1&z&(z^{-1},z^{-1})\\
z&1&1&z^{-1}&(z^2,z)\\
1&1&z&1&(1,z).
\end{array}
\]
The first two rows prove the homomorphism condition. Suppose the third row occurs. The pair $(1,z)\in\Gamma$ and exterior phase $z$ give $1+2z$ in the spectrum of $T$. Applying Lemma~\ref{forced} again to a matching eigenvalue of $T^+$ gives the possibilities
\[
\begin{array}{c|c|c|c|c}
A&B&C&z'&(u,v)\\\hline
z&z&1&1&(z,1)\\
z&1&z&z^{-1}&(z^2,z^2)\\
1&z&z&z&(z^{-1},1).
\end{array}
\]
In the first row, multiply $(z,1)$ by $(1,z)$ to obtain $(z,z)$ in $\Gamma$. In the second row, the diagonal pair $(z^2,z^2)$ generates the same subgroup as $(z,z)$, since $2$ is invertible modulo $p$. In the third row, raise $(z^{-1},1)$ to a suitable power to obtain $(z^2,1)$, and multiply by $(1,z)$ to obtain $(z^2,z)$. Each case gives one of the pairs in Lemma~\ref{dualh}. Necessity and all equivalences follow.
\end{proof}

\section{The ternary spectra and the beta condition}
Fix $\omega=e^{2\pi i/3}$. For $x=(u,v)\in\Gamma$, write
\begin{equation}\label{LR}
\begin{aligned}
L_1(x)&=v-\omega^2u,& L_2(x)&=v-\omega u,\\
R_1(x)&=-u+\omega v,& R_2(x)&=-u+\omega^2v.
\end{aligned}
\end{equation}
These are the two nontrivial phase sectors for $T$ and $T^+$. The trivial sector $2u+v$ is the same on both sides. We may cancel this common multiset, even when its values also occur in the other sectors. Thus cospectrality is equivalent to equality of the multisets of all $L_1,L_2$ and all $R_1,R_2$.

\begin{lemma}\label{betasuff}
If $\beta\Gamma=\Gamma$, then the two digraphs are cospectral.
\end{lemma}
\begin{proof}
Direct substitution gives $L_1(x)=R_2(\beta x)$ and $L_2(x)=R_1(\beta x)$. Since $\beta^2=\mathrm{id}$, these identities provide a bijection of the nontrivial occurrences, preserving multiplicities. The trivial occurrences already agree.
\end{proof}

\begin{proof}[Proof of Theorem~\ref{main3}, part 1]
Set $e=\Exp(H)$ and $F=\Q(\mu_e)$. Since $3\nmid e$, the standard degree formula for cyclotomic fields gives $[F(\omega):F]=\varphi(3e)/\varphi(e)=2$. Hence $1,\omega$ are linearly independent over $F$. In that basis,
\[
L_1=(u+v)+u\omega,\quad L_2=v-u\omega,\quad
R_1=-u+v\omega,\quad R_2=-(u+v)-v\omega.
\]
Every coefficient $u,v$ is nonzero. No nontrivial-sector value lies in $F$, whereas every trivial-sector value does. A matching $L_1(u,v)=R_2(u',v')$ forces $v'=-u$ and $u'=-v$. A matching $L_2(u,v)=R_1(u',v')$ gives the same pair $\beta(u,v)$.

A same-sector matching $L_1=R_1$ would force $u'=-u-v$, so $|u+v|=1$. For unit complex numbers, this implies $u/v\in\{\omega,\omega^2\}$, by expanding the squared modulus. This is impossible because $u/v\in\mu_e$ and $3\nmid e$. The matching $L_2=R_2$ leads to the same impossibility. Every occurrence of $L_1(x)$ must therefore match $R_2(\beta x)$ with $\beta x\in\Gamma$. We obtain $\beta\Gamma\subseteq\Gamma$, and the involution property gives equality. Sufficiency is Lemma~\ref{betasuff}.
\end{proof}

\section{Necessity in the presence of three torsion}
\begin{lemma}[Unit-circle chords]\label{chord}
Let $a,b,c,d$ have modulus one and $a\ne b$. If $a-b=c-d$, then $(c,d)=(a,b)$ or $(c,d)=(-b,-a)$. The two possibilities may coincide when $a=-b$.
\end{lemma}
\begin{proof}
For $\delta=a-b\ne0$, the possible endpoints $d$ satisfy $|d|=|d+\delta|=1$. Two distinct circles have at most two intersections. The known solutions $b$ and $-a$ exhaust these possibilities; if they coincide, the circles are tangent.
\end{proof}

\begin{lemma}\label{beta3h}
If $3\mid|H|$ and $\beta\Gamma=\Gamma$, then the homomorphism condition holds.
\end{lemma}
\begin{proof}
Put $B=(-1,-1)$ and $s(u,v)=(v,u)$, so $\beta x=B s(x)$. As $\beta(1,1)=B$, invariance gives $B\in\Gamma$ and $s\Gamma=\Gamma$. By finite character duality, exchange of $\kappa,\kappa'$ induces an automorphism of $H$; alternatively use Proposition~\ref{lattice}. It induces an automorphism of $V=H/3H$, a nonzero vector space over $\Fthree$ of dimension at most two.

If $\dim V=2$, the classes of the two generators form a basis, and the linear functional taking value $1$ on both gives $h$. If $\dim V=1$, exchange invariance excludes exactly one generator class being zero. Since the two classes generate $V$, both are nonzero. Normalize a nonzero functional to take value $1$ on $\kappa'$; its value on $\kappa$ is $1$ or $2$. Composing with $H\to V$ completes the proof.
\end{proof}

\begin{proof}[Proof of Theorem~\ref{main3}, part 2]
Assume cospectrality and suppose, for contradiction, that the homomorphism condition is absent. By Lemma~\ref{dualh}, none of $(\omega,\omega)$, $(\omega^2,\omega^2)$, $(\omega^2,\omega)$ and $(\omega,\omega^2)$ belongs to $\Gamma$.

Apply Lemma~\ref{chord} to a matching for $L_1(1,1)=1-\omega^2\ne0$. The possible target pairs are
\[
(\omega^2,\omega^2),\quad (\omega^2,\omega),\quad
B=(-1,-1),\quad A=(-1,-\omega).
\]
The first two are excluded. If $A\in\Gamma$, apply the chord lemma to $L_1(A)=\omega^2-\omega$. Its targets are
\[
(-\omega^2,-1),\quad (\omega,\omega),\quad
(-\omega^2,-\omega^2),\quad (\omega,1).
\]
The second gives a forbidden diagonal pair; the square of the third does likewise. The square of the first is $(\omega,1)$. Since $A^2=(1,\omega^2)$, the first and fourth cases give $(\omega,\omega^2)$ by multiplication with $A^2$. Thus $A\notin\Gamma$ and necessarily $B\in\Gamma$.

Write $a=(1,\omega)$. For $x=(u,v)$ with $L_1(x)\ne0$, the chord lemma gives
\[
\begin{array}{c|cc}
\text{matching}&\text{same chord}&\text{opposite chord}\\\hline
L_1(x)=R_1(y)&y=(\omega^2u,\omega^2v)&y=a\beta x\\
L_1(x)=R_2(y)&y=(\omega^2u,\omega v)&y=\beta x.
\end{array}
\]
If a same-chord target lay in $\Gamma$, division by $x\in\Gamma$ would yield a forbidden pair. Consequently
\begin{equation}\label{constraints}
\begin{split}
L_1(x)\ne0&\ \Longrightarrow\ \beta x\in\Gamma\text{ or }a\beta x\in\Gamma,\\
L_2(x)\ne0&\ \Longrightarrow\ \beta x\in\Gamma\text{ or }a^2\beta x\in\Gamma.
\end{split}
\end{equation}
The second implication follows by the same calculation: its same-chord multipliers are $(\omega,\omega)$ and $(\omega,\omega^2)$, also forbidden.

If $\beta\Gamma=\Gamma$, Lemma~\ref{beta3h} contradicts our assumption. We may therefore assume $\beta\Gamma\ne\Gamma$. Since $B\in\Gamma$ and $s(B)=B$, the set $G_1=s\Gamma=\beta\Gamma$ is a subgroup. Put $K=G_1\cap\Gamma$ and $q(y_1,y_2)=y_2/y_1$. Also $a\notin\Gamma$, since $Ba=A\notin\Gamma$.

For $y=\beta x\in G_1\setminus K$, if $q(y)\notin\{\omega,\omega^2\}$, both $L_1(x)$ and $L_2(x)$ are nonzero. The two implications~\eqref{constraints} give $ay,a^2y\in\Gamma$; their quotient gives $a\in\Gamma$, a contradiction. Therefore every $y\in G_1\setminus K$ has $q(y)\in\{\omega,\omega^2\}$.

Fix such a $y$. For every $k\in K$, $yk\notin K$, so $q(y)q(k)\in\{\omega,\omega^2\}$. The subgroup $q(K)$ is contained either in $\{1,\omega\}$ or in $\{1,\omega^2\}$. Neither two-element set is a subgroup, so $q(K)=\{1\}$. It follows that
\[
K=\{(d,d):d\in D\},\qquad \ker(q|_{G_1})=K,\qquad q(G_1)=\mu_3.
\]
The last equality uses that $G_1\ne K$ and $q(G_1)$ is a subgroup. Thus $[G_1:K]=[\Gamma:K]=3$. Since $s\Gamma=G_1$, the image $q(\Gamma)$ is also $\mu_3$. Choose $x_0=(\alpha,\omega\alpha)\in\Gamma$. We have
\[
\Gamma=K\sqcup x_0K\sqcup x_0^2K,
\]
and $-1\in D$ because $B\in K$.

For $x\in K$, the combined $L$ and $R$ multisets agree: on a diagonal pair, $R_2=-L_1$ and $R_1=-L_2$, and multiplication by $-1$ permutes $D$. Cancel them. Let $s_0=\omega-\omega^2$. On the remaining cosets, the nonzero values are
\[
\begin{array}{c|cc}
\text{coset}&T&T^+\\\hline
x_0K&D\alpha s_0&D\omega\alpha s_0\\
x_0^2K&D\alpha^2 s_0&D\omega^2\alpha^2s_0.
\end{array}
\]
Here signs have been absorbed in $D$; each coset contributes $|D|$ zero values as well. Distinct cosets of $D$ are disjoint. Equality with multiplicities therefore implies
\[
\{D\alpha,D\alpha^2\}=\{D\omega\alpha,D\omega^2\alpha^2\}
\]
as multisets of two cosets. The direct permutation gives $\omega\in D$, hence $(\omega,\omega)\in\Gamma$. The crossed permutation gives $\omega^2\alpha\in D$; dividing $x_0$ by the corresponding diagonal pair produces $(\omega,\omega^2)\in\Gamma$. Both contradict absence of $h$. Necessity is proved. Lemma~\ref{suffh} gives sufficiency and the isomorphism statements.
\end{proof}

\section{Relation lattices and consequences}
Let
\[
L=\ker\bigl[\Z^2\to H,\ (r,t)\mapsto r\kappa+t\kappa'\bigr].
\]
The letter $t$ here is an integer coordinate, not a support element.
\begin{proposition}\label{lattice}
The homomorphism condition is equivalent to the existence of $b\in\{1,2\}$ such that
\[
L\subseteq\{(r,t):br+t\equiv0\pmod p\}.
\]
The condition $\beta\Gamma=\Gamma$ is equivalent to both
\[
s(L)=L\quad\text{and}\quad r+t\equiv0\pmod2\text{ for every }(r,t)\in L,
\]
where $s(r,t)=(t,r)$.
\end{proposition}
\begin{proof}
A functional $(r,t)\mapsto br+t$ descends from $\Z^2$ to $H$ precisely when it vanishes on $L$. For the second assertion, $\Gamma$ is the annihilator of $L$: its elements satisfy $u^rv^t=1$ for all $(r,t)\in L$. Therefore $s\Gamma=\Gamma$ is equivalent to $sL=L$ by finite duality, and $(-1,-1)\in\Gamma$ is equivalent to evenness of every relation sum. Since $\beta x=(-1,-1)s(x)$, subgroup invariance is equivalent to these two conditions.
\end{proof}

\begin{corollary}
If $p\geq5$ and $p\nmid|H|$, the pair is not cospectral. If the homomorphism condition holds for any odd $p$, then $p\mid|H|$ and $H\times\Fp$ is not cyclic.
\end{corollary}
\begin{proof}
The prescribed value $h(\kappa')=1$ makes $h$ surjective. Hence $p\mid|H|$, and $H$ has nontrivial $p$-torsion. The $p$-torsion subgroup of $H\times\Fp$ has rank at least two. A cyclic group has $p$-torsion rank at most one.
\end{proof}

\begin{corollary}\label{cyclicbeta}
Suppose $H=C_n$ and $\beta\Gamma=\Gamma$. Then both $\kappa$ and $\kappa'$ generate $H$. After choosing coordinates $\kappa=1$ and $\kappa'=m$, the beta condition is equivalent to
\[
n\text{ even},\qquad \gcd(m,n)=1,\qquad m^2\equiv1\pmod n.
\]
Distinctness additionally requires $m\not\equiv1\pmod n$.
\end{corollary}
\begin{proof}
Exchange invariance gives an automorphism interchanging the two generators, so they have equal order. In a cyclic group they generate the same subgroup, which must be all of $H$. Exchange is multiplication by $m$, and exchanging back requires $m^2=1$ modulo $n$. A character taking value $-1$ on $1$ exists precisely when $n$ is even; it also takes value $-1$ on $m$ because every unit modulo an even integer is odd. This proves both directions.
\end{proof}

\section{Examples and comparison with earlier work}
\begin{example}[A previously studied pair on twelve vertices]\label{twelve}
Let $p=3$, $H=C_4$, $\kappa=1$ and $\kappa'=3$. Exchange is multiplication by $3$, and the character $x\mapsto(-1)^x$ gives the parity condition. Thus $\beta\Gamma=\Gamma$, and Theorem~\ref{main3} gives cospectrality.
The group isomorphism $F:C_4\times C_3\to C_{12}$, $F(x,y)=9x+4y$, gives
\[
F(T)=\{1,3,9\},\qquad F(T^+)=\{1,5,7\}.
\]
Brown's Example~3.1~\cite{Brown} has supports $\{1,5,11\}$ and $\{1,3,9\}$; multiplication by $5$ sends the first to $F(T^+)$. M\"onius's Example~3.1.13~\cite{Moenius} has supports $\{5,7,11\}$ and $\{3,9,11\}$, obtained by applying the same multiplier $11$ to $F(T^+)$ and $F(T)$ respectively. This pair is therefore an antecedent example. The published examples also establish nonisomorphism for this case.
\end{example}

\paragraph{Overlap with the 2020 article.}
Theorem~10 of M\"onius~\cite[pp.~54--55]{Moenius2020} recovers adjacency cospectrality for an entire cyclic subfamily, beyond the twelve-vertex example: take $p=3$, $H=C_k$, $3\nmid k$, $4\mid k$, $\kappa=1$ and $\kappa'=\mu=1+k/2$. Put $n=3k$, and let $F:C_k\times C_3\to C_n$ be the CRT isomorphism determined by $F(x,y)\equiv x\pmod{k}$ and $F(x,y)\equiv y\pmod{3}$. Define
\[
S=F(T),\qquad P=F(T^+),\qquad b=F(\mu,2),\qquad a_0=F(1,0),
\]
and the odd equilibrium progression of length six
\[
\mathrm{Eq}=\{a_0+j n/6\pmod n:0\leq j<6\}.
\]
Here $n\in4\mathbf N$, $n/6=k/2$ is even, $b$ is a unit modulo $n$, $S$ consists of exactly half the elements of $\mathrm{Eq}$, and $bP=\mathrm{Eq}\setminus S$. The complement branch of Theorem~10, with the additional set empty, gives cospectrality of $S$ and $bP$; multiplication by the unit $b$ preserves the adjacency spectrum. Thus it also gives cospectrality of $S$ and $P$ for every $k$ in this subfamily. For $k=4$, these supports are those in Example~\ref{twelve}, and $b=11$.

This correspondence is a deduction from the inspected primary article, not a quotation or a claim that the article names this family. The comparison concerns cospectrality only: it establishes no antecedence conclusion for spectral injectivity (source witness 10B1), a nonisomorphism criterion, or the general $h$/$\beta$ classifications. The priority status remains explicitly \textbf{N = NON AUDITEE} for those claims, the other involutions, and the complete cyclic family. Primary-source comparison with Meng (1998) and Huang--Chang (2001), the complete literature audit, and independent human mathematical review remain open.

\begin{example}[Homomorphism and obstruction]
For $H=C_p$ with $\kappa=2$, $\kappa'=1$, the identity homomorphism gives the reflection isomorphism for every odd $p$. A negative example takes $p=5$, $H=C_5$, $\kappa=0$ and $\kappa'=1$. Every homomorphism sends $0$ to $0$, so the pair is not cospectral. This example retains the loop in $T$ under our convention. For a loop-free obstruction take $\kappa=3$, $\kappa'=1$ in $C_5$.
\end{example}

\begin{example}[Noncyclic positive regime]
Let $H=C_p\times C_p$, $\kappa=(1,0)$ and $\kappa'=(0,1)$. The functional $h(x_1,x_2)=x_1+x_2$ gives a positive case for every odd $p$, including primes at least five. The ambient group is $C_p^3$ and is not a circulant group.
\end{example}

Several distinctions are relevant to priority. Brown's cospectral construction on order $2^r q$ has valence $q$ for an odd prime $q$; its ternary instance can meet our family, whereas $q\geq5$ does not preserve the three-element support. The circulant criteria of Mans et al.~\cite{MPS} and M\"onius~\cite{Moenius} have their own arithmetic or support hypotheses. They are not automatically statements about a general noncyclic $H\times C_p$. The correction by Qin et al.~\cite{Qin} also illustrates the need to retain the exact hypotheses of spectrum-determination results.

Wiedemann and Zieve~\cite[Theorem~1.1]{WZ} give equivalences for sparse circulants involving affine support maps, independent row and column permutations, and autocorrelation matrices $AA^{\mathsf T}$. This invariant must be distinguished from adjacency spectrum. Indeed, translation of any support multiplies its adjacency matrix by a permutation matrix $P$, up to the choice of row convention, and $AP(AP)^*=AA^*$. This holds for every support translation, without either classification criterion. Finally, CI statements, such as those of Dobson and Morris~\cite{DM}, concern isomorphisms and group automorphisms rather than an implication from cospectrality to isomorphism.

\section{Exact computational checks}
The accompanying script represents $H=\Z^2/L$ using bases $(a,0),(b,d)$, with $a,d>0$ and $0\leq b<a$. It excludes coincident generator classes. If $N=pad$, a character has values $\xi^U,\xi^V$ where $\xi$ is primitive of order $N$ and
\[
U=\frac{N i}{a},\qquad
V=p(ja-bi)\pmod N,\qquad 0\leq i<a,\quad 0\leq j<d.
\]
This follows from $u^a=1$ and $u^bv^d=1$. Exterior phases are $\xi^{k ad}$. Each eigenvalue is stored as an integer coefficient vector in $\Z[X]/(\Phi_N(X))$; multiset comparison is therefore exact, with no floating-point tolerance. The predicates $h$ and $\beta$ are computed independently from the relation lattice, not inferred from the spectrum. A second check forms adjacency matrices for selected small examples and compares their characteristic polynomials.

\begin{center}
\begin{tabular}{rrrrr}
\toprule
$p$ & Bound on $|H|$ & Presentations & Cospectral & Mismatches\\
\midrule
3 & 120 & 11853 & 2668 & 0\\
5 & 30 & 732 & 59 & 0\\
7 & 20 & 319 & 6 & 0\\
11 & 12 & 115 & 1 & 0\\
\bottomrule
\end{tabular}
\end{center}

These checks are finite validations of the implementation and examples. They do not replace the proofs. For $p=3$ through index 120, the separately executed code recovers the two regime counts reported by the source notes: 165 cospectral presentations among 6948 without three torsion, and 2503 among 4905 with three torsion. The other historical scans are not reproduced. The script, its version requirements and its executed output are included with this version.

\section*{Preparation and responsibility}
This review version was prepared with generative AI assistance for proof exposition, literature comparison and code generation. No independent human review is certified in this version. Before submission, the named author must check each argument, approve the wording and authorship, and confirm the literature positioning. The technical publication coordinator is not assigned mathematical authorship by virtue of that role. Funding, competing interests and the final assistance declaration require the author's confirmation.

\begin{thebibliography}{99}
\bibitem{PR} B.~Poonen and M.~Rubinstein, \emph{The Number of Intersection Points Made by the Diagonals of a Regular Polygon}, SIAM J. Discrete Math. \textbf{11} (1998), 133--156. \href{https://doi.org/10.1137/S0895480195281246}{doi:10.1137/S0895480195281246}.
\bibitem{Brown} J.~Brown, \emph{Isomorphic and Nonisomorphic, Isospectral Circulant Graphs}, arXiv:0904.1968 (2009). \url{https://arxiv.org/abs/0904.1968}.
\bibitem{Moenius2020} K.~M\"onius, \emph{Constructions of isospectral circulant graphs}, Elemente der Mathematik \textbf{75}(2) (2020), 45--57. \href{https://doi.org/10.4171/EM/404}{doi:10.4171/EM/404}.
\bibitem{Moenius} K.~M\"onius, \emph{Algebraic and Arithmetic Properties of Graph Spectra}, doctoral thesis, Universit\"at W\"urzburg (2021). \href{https://doi.org/10.25972/OPUS-23085}{doi:10.25972/OPUS-23085}.
\bibitem{MPS} B.~Mans, F.~Pappalardi and I.~E.~Shparlinski, \emph{On the spectral \`Ad\'am property for circulant graphs}, Discrete Math. \textbf{254} (2002), 309--329. \href{https://doi.org/10.1016/S0012-365X(01)00374-0}{doi:10.1016/S0012-365X(01)00374-0}.
\bibitem{LL} T.~Y.~Lam and K.~H.~Leung, \emph{On Vanishing Sums of Roots of Unity}, J. Algebra \textbf{224} (2000), 91--109. \href{https://doi.org/10.1006/jabr.1999.8089}{doi:10.1006/jabr.1999.8089}.
\bibitem{WZ} D.~Wiedemann and M.~E.~Zieve, \emph{Equivalence of Sparse Circulants: The Bipartite \`Ad\'am Problem}, arXiv:0706.1567 (2007). \url{https://arxiv.org/abs/0706.1567}.
\bibitem{DM} E.~Dobson and J.~Morris, \emph{Toida's Conjecture is True}, Electronic J. Combin. \textbf{9}(1) (2002), R35. \url{https://www.cs.uleth.ca/~morris/Research/Toida.pdf}.
\bibitem{LZ} X.~Liu and S.~Zhou, \emph{Eigenvalues of Cayley graphs}, Electronic J. Combin. \textbf{29}(2) (2022), P2.9. \href{https://doi.org/10.37236/8569}{doi:10.37236/8569}.
\bibitem{Qin} Z.~Qin, X.~Xu, J.~Meng and G.~Wang, \emph{Correction to ``A Classification of Spectrum-determined Circulant Digraphs''}, Acta Math. Appl. Sinica, English Series \textbf{37}(4) (2021), 867--868. \href{https://doi.org/10.1007/s10255-021-1052-6}{doi:10.1007/s10255-021-1052-6}.
\end{thebibliography}
\end{document}
